\documentclass[11pt]{article}
\usepackage[margin=1in]{geometry}
\usepackage{booktabs}
\usepackage{longtable}
\usepackage{enumitem}
\usepackage{amsmath}
\usepackage{hyperref}
\usepackage{xcolor}
\usepackage{fancyvrb}

\hypersetup{colorlinks=true,linkcolor=blue,urlcolor=blue}

\title{Assessment of the \texttt{carryover\_focus/} Directory}
\author{pmsimstats2025 Project Notes}
\date{March 2025}

\begin{document}
\maketitle

\section{Summary}

This directory is a self-contained sub-analysis that distills the
144-condition full simulation (\texttt{simulation\_clustered.R}) down to
9~conditions, isolating a single question: what happens to statistical
power for biomarker moderation detection when pharmacological carryover
is present but the analysis model ignores it?

\section{Contents}

\begin{table}[h]
\centering
\begin{tabular}{@{}lp{10cm}@{}}
\toprule
File & Role \\
\midrule
\texttt{PLAN.md}
  & Design document with line-level extraction references to the
    parent script \\
\texttt{README.md}
  & User-facing documentation (how to run, expected results) \\
\texttt{carryover\_power\_simulation.R}
  & 701-line self-contained simulation (functions extracted verbatim
    from \texttt{simulation\_clustered.R}) \\
\texttt{predictor\_comparison.tex/.pdf}
  & 7-page technical note documenting a key finding that emerged
    during this work \\
\texttt{output/}
  & Three \texttt{.RData} result sets and corresponding figures for
    different \texttt{tsd\_min} values \\
\bottomrule
\end{tabular}
\end{table}

\section{Intellectual Arc}

The work progressed through three distinct phases.

\subsection{Phase 1: Simplification}

The full parameter grid in \texttt{simulation\_clustered.R} sweeps
4~designs $\times$ 2~sample sizes $\times$ 4~biomarker moderation
levels $\times$ 3~carryover values $\times$ 2~model specifications
$= 144$ conditions. This sub-analysis reduces the grid to carryover
half-life $\times$ model specification for the Hybrid design alone.

A new predictor, \texttt{naive\_drug\_weeks}, was introduced as the
``ignores carryover'' comparator. It replaces the original
\texttt{weeks\_on\_drug}, which was structurally deficient: because
\texttt{weeks\_on\_drug} $=$ \texttt{cumsum(treatment)} never
decreases during off-drug periods, it provides no on/off contrast and
yields power $\approx 0.10$ regardless of the carryover level. The
replacement \texttt{naive\_drug\_weeks} $=$
\texttt{if\_else(treatment == 1, weeks\_on\_drug, 0)} drops to zero
when off drug, preserving the cumulative scale while restoring the
contrast needed for the interaction test.

\subsection{Phase 2: Discovery of the $t_{\text{sd}} = 0$
  Discontinuity}

The initial results (the \texttt{tsd0} output set) showed a
step-function pattern: power drops from $\sim$0.81 at $t_{1/2} = 0$
to $\sim$0.20 at $t_{1/2} = 0.05$ weeks (8.4~hours), then remains
flat regardless of further increases in half-life.

The \texttt{predictor\_comparison.tex} document traces this to a
mathematical artifact. At the exact moment of discontinuation,
$t_{\text{sd}} = 0$, and the carryover function evaluates as:
\[
  \left(\frac{1}{2}\right)^{0 / t_{1/2}}
  = \left(\frac{1}{2}\right)^{0} = 1
  \quad \text{for any } t_{1/2} > 0
\]
Since weeks~10 and~20 in the Hybrid design have $t_{\text{sd}} = 0$,
they carry full drug contamination regardless of whether the half-life
is 8~hours or 34~hours. These two timepoints constitute half of the
off-drug measurements, producing severe and constant attenuation of
the on/off contrast. This same artifact is present in Hendrickson
et~al.'s original \texttt{pmsimstats} code.

\subsection{Phase 3: The $t_{\text{sd,min}}$ Correction}

Setting \texttt{tsd\_min}~$= 0.5$~weeks imposes a realistic gap
between the last dose and the first off-drug measurement, restoring
continuity in the carryover function:
\[
  t_{\text{sd,eff}} = \max(t_{\text{sd}},\; t_{\text{sd,min}})
\]
The corrected simulation (the \texttt{tsd05} output set) produces the
expected gradual, monotonic power decline as the true carryover
half-life increases.

\section{The Three Output Sets}

\begin{table}[h]
\centering
\caption{Output sets generated by \texttt{carryover\_power\_simulation.R}
  under different $t_{\text{sd,min}}$ values.}
\label{tab:outputs}
\begin{tabular}{@{}llp{8cm}@{}}
\toprule
Suffix & $t_{\text{sd,min}}$ & Pattern \\
\midrule
\texttt{\_tsd0}
  & 0 (default)
  & Step-function cliff. Modeled power $\sim$0.80 across all
    half-lives; naive power drops to $\sim$0.20 at any nonzero
    half-life and stays flat. This is the discontinuity artifact. \\[6pt]
\texttt{\_tsd05}
  & 0.5 weeks
  & Gradual decline. Both lines start at $\sim$0.80 at carryover
    $= 0$; the naive line declines smoothly as the half-life
    increases while the modeled line holds steady. This is the
    pharmacologically expected pattern. \\[6pt]
(no suffix)
  & 0.5 weeks
  & Appears identical to the \texttt{tsd05} set based on the
    figure, suggesting the most recent run used
    $t_{\text{sd,min}} = 0.5$. \\
\bottomrule
\end{tabular}
\end{table}

\section{Relationship to the Broader Project}

This sub-analysis serves three purposes for the paper:

\begin{enumerate}
  \item \textbf{Pedagogical figure.} The two-panel power/effect-size
    line plot is far more readable than the full simulation's heatmaps
    for communicating the carryover finding to a clinical audience.

  \item \textbf{Identifies an artifact in Hendrickson et~al.\ (2020).}
    The $t_{\text{sd}} = 0$ discontinuity explains why the published
    simulation heatmaps show an implausibly sharp power cliff at tiny
    carryover half-lives. The artifact affects all trial designs that
    place measurements at the treatment boundary.

  \item \textbf{Proposes a correction.} The $t_{\text{sd,min}}$
    parameter eliminates the mathematical discontinuity and produces
    pharmacologically realistic power curves. The recommendation is to
    report both the uncorrected (sharp cliff) and corrected (gradual
    decline) results, demonstrating that the power-cliff finding is
    partly an artifact of the $t_{\text{sd}} = 0$ assumption.
\end{enumerate}

\section{Open Items}

\begin{itemize}
  \item Per \texttt{PLAN.md}, the ``run and verify'' and ``compare
    against \texttt{simulation\_clustered\_results.RData}'' steps
    remain unchecked, although the output artifacts exist. Formal
    verification against the parent simulation's results would confirm
    numerical consistency between the simplified and full codebases.

  \item The \texttt{PLAN.md} checklist has not been updated to reflect
    the completed work (script creation, README, output generation).
\end{itemize}

\bigskip
\hrule
\medskip
{\small\itshape
  Rendered on 2026-03-04 at 09:12 PST.\\
  Source: \texttt{\~/prj/alz/01-pmsimstats/pmsimstats2025/%
    analysis/scripts/carryover\_focus/directory\_assessment.tex}
}

\end{document}
